\documentclass[12pt]{extarticle}
\def\activityid{F03-U-v3}\usepackage[letterpaper,margin=.65in,top=.60in,bottom=.65in]{geometry}
\usepackage{fix-cm}
\usepackage[T1]{fontenc}\usepackage[default]{sourcesanspro}
\usepackage{microtype,tikz,array,tabularx,fancyhdr,amsmath,amssymb}
\usepackage[hidelinks]{hyperref}\usetikzlibrary{calc,arrows.meta}
\setlength{\parindent}{0pt}\setlength{\parskip}{7pt}\setlength{\headheight}{0pt}\setlength{\footskip}{26pt}
\setlength{\emergencystretch}{2em}\pagestyle{fancy}\fancyhf{}\renewcommand{\headrulewidth}{0pt}\renewcommand{\footrulewidth}{.4pt}
\fancyfoot[L]{\scriptsize Bellingham Math Circle / Week 3 / \activityid}\fancyfoot[R]{\scriptsize \thepage}
\definecolor{ink}{gray}{.12}\definecolor{soft}{gray}{.95}\color{ink}
\newcommand{\PageTitle}[2]{{\fontsize{10}{12}\selectfont\bfseries WEEK 3\quad CODE MACHINES\hfill #1\par}\vspace{3pt}{\fontsize{26}{29}\selectfont\bfseries #2\par}\vspace{2pt}}
\newcommand{\NameLine}{{\small Name: \rule{2.65in}{.4pt}\hfill Date: \rule{1.35in}{.4pt}\par}\vspace{4pt}}
\newcommand{\Task}[2]{\par\vspace{3pt}{\large\bfseries #1}\par #2\par}
\newcommand{\Subhead}[1]{\par\vspace{3pt}{\large\bfseries #1}\par}
\newcommand{\RuleBox}[1]{\par\begingroup\setlength{\fboxsep}{8pt}\colorbox{soft}{\parbox{\dimexpr\linewidth-16pt\relax}{#1}}\endgroup\par}
\newcommand{\WriteBox}[2]{\par\textbf{#1}\par\vspace{2pt}\begin{tikzpicture}\draw[black!40,line width=.5pt] (0,0) rectangle (\dimexpr\linewidth-.5pt\relax,#2);\end{tikzpicture}\par}
\newcommand{\StudentFooter}{\vfill{\small Try it with objects. Explain by pointing, talking, or drawing.}\par}
\newcommand{\LampRing}[3]{\begin{tikzpicture}[x=1in,y=1in]
\foreach \i in {1,...,#1}{\coordinate (v\i) at ({90-360*(\i-1)/#1}:#2);}
\foreach \i in {1,...,#1}{\pgfmathtruncatemacro{\j}{mod(\i,#1)+1}\draw[thick] (v\i)--node[midway,fill=white,inner sep=2pt,font=\small]{\i\j}(v\j);}
\foreach \i in {1,...,#1}{\node[circle,draw,fill=white,minimum size=.52in,inner sep=0pt,font=\bfseries] at(v\i){\i};}
\foreach \i in {#3}{\node[circle,draw,fill=ink,text=white,minimum size=.52in,inner sep=0pt,font=\bfseries] at(v\i){\i};}
\end{tikzpicture}}
\newcommand{\Lights}[1]{\begin{tikzpicture}[x=.55in,y=.55in,baseline=-.10in]\foreach \v [count=\i] in {#1}{\ifnum\v=1\fill (\i,0) circle(.16in);\else\draw (\i,0) circle(.16in);\fi}\end{tikzpicture}}
\newcommand{\ClockRing}[2]{\begin{tikzpicture}[x=1in,y=1in]\draw[black!30] (0,0) circle(#2);\pgfmathtruncatemacro{\n}{#1-1}\foreach\i in {0,...,\n}{\node[circle,draw,fill=white,minimum size=.35in,inner sep=0pt] at({90-360*\i/#1}:#2){\i};}\draw[-{Stealth},thick] (-.35,.15) arc(155:25:.38);\node[font=\small] at(0,-.18){clockwise};\end{tikzpicture}}
\newcommand{\Key}[2]{\begin{center}\renewcommand{\arraystretch}{1.55}\begin{tabular}{c|*{#1}{c}}in&#2\end{tabular}\end{center}}


% v3 student pages use one running header and numbered tasks only.
\newcommand{\StudentHeader}[1]{{\fontsize{10}{12}\selectfont\bfseries Week 3 / Code machines / #1\par}\vspace{10pt}}
\newcommand{\Problem}[2]{\par\vspace{4pt}\textbf{Problem #1:} #2\par}
\newcommand{\Space}[1]{\begin{tikzpicture}\draw[black!35] (0,0) rectangle (\dimexpr\linewidth-.5pt\relax,#1);\end{tikzpicture}\par}
\newcommand{\Record}[2]{#1\par\Space{#2}}

\begin{document}
\StudentHeader{Grades 4--5}

\Problem{1}{Use five letter cards. One turn changes every letter using the chosen key, keeping message positions fixed. Start each experiment with ABCDE. Repeat the same key and record each message until the whole row first returns. Circle that turn.}
\begin{center}\renewcommand{\arraystretch}{1.6}\begin{tabular}{c|ccccc}Input&A&B&C&D&E\\\hline Key R&B&C&D&E&A\\Key S&B&C&A&E&D\end{tabular}\end{center}
\begin{center}\renewcommand{\arraystretch}{1.9}\begin{tabular}{c|p{2.3in}|p{2.3in}}Turn&R message&S message\\\hline0&ABCDE&ABCDE\\1&&\\2&&\\3&&\\4&&\\5&&\\6&&\\\end{tabular}\end{center}
\Problem{2}{For each key, follow A and draw its arrows until it comes back. Then start at a letter not yet drawn. Draw every loop, including any letter that stays put. Compare the two keys' loop sizes and return times.}
\Space{1.35in}

\newpage
\StudentHeader{Grades 4--5}

\Problem{3}{Build three more keys from five letter cards: loops of sizes 4+1, then 2+2+1, then 3+1+1. Draw each key and test when all letters first return.}
\begin{center}\renewcommand{\arraystretch}{3}\begin{tabular}{l|p{3.1in}|p{.85in}}Loop sizes&My arrows&Return\\\hline4 + 1&&\\2 + 2 + 1&&\\3 + 1 + 1&&\\\end{tabular}\end{center}
\Problem{4}{Return to key S, with loops of sizes 3 and 2. Move a counter around each loop, one arrow per turn. In each row below, circle every turn when that counter is home. Box the first positive turn circled in both rows.}
\begin{center}\renewcommand{\arraystretch}{2}\begin{tabular}{c|rrrrrrrrrrrr}Loop&1&2&3&4&5&6&7&8&9&10&11&12\\\hline Size 3&1&2&3&4&5&6&7&8&9&10&11&12\\Size 2&1&2&3&4&5&6&7&8&9&10&11&12\end{tabular}\end{center}
\Problem{5}{Invent another five-letter key. Predict its return from its loops, then test. Record your prediction and result.}
\Space{1.2in}

\newpage
\StudentHeader{Grades 4--5}

\Problem{6}{Arrange five letter cards into loops. Organize different arrangements by their largest loop: 5, then 4, then 3, then 2, then 1. For each largest loop, split the remaining cards in every way you can. Record sizes and the first common return. Changing letter names does not create a new list of sizes.}
\begin{center}\renewcommand{\arraystretch}{2.2}\begin{tabular}{p{2.5in}|p{2.5in}}Loop sizes&First common return\\\hline&\\&\\&\\&\\&\\&\\&\\&\\\end{tabular}\end{center}
\Problem{7}{Choose the arrangement with the longest return. Write a key that has those loops. Test it with counters and mark its return turns. Compare with your earlier keys.}
\Space{1.4in}

\newpage
\StudentHeader{Grades 4--5}

\Problem{8}{Key P swaps A and B. Key Q swaps B and C. Both leave all other letters alone. Start with ABCDE each time. Apply the first key, record the whole row, then apply the second key to that row.}
\begin{center}\renewcommand{\arraystretch}{2.3}\begin{tabular}{l|p{1.8in}|p{1.8in}}Order&After first key&After second key\\\hline P then Q&&\\Q then P&&\\\end{tabular}\end{center}
\Problem{9}{Treat P then Q as one turn. Repeat that whole turn on ABCDE until the row first returns. Draw the loops of this combined key.}
\Space{1.3in}
\Problem{10}{Change Q to a new key T that swaps D and E and fixes A,B,C. Test P then T and T then P. Compare this pair with Problem 8.}
\Space{1.1in}
\Problem{11}{Start with ABCDE and apply P then Q. Now find an order of these same two keys that undoes the result. Record the full check.}
\Space{.85in}

\end{document}
